\documentclass[12pt,a4paper]{article}

\usepackage[utf8]{inputenc}
\usepackage[T1]{fontenc}
\usepackage[margin=2.5cm]{geometry}
\usepackage{hyperref}
\usepackage{microtype}
\usepackage{parskip}
\usepackage{lmodern}

\hypersetup{colorlinks=true, urlcolor=blue!60!black}

\pagestyle{empty}

\begin{document}

\begin{flushright}
Kundai Farai Sachikonye \\
Auf der Lichtung 19, 82178 Puchheim \\
\href{https://github.com/fullscreen-triangle}{github.com/fullscreen-triangle} \\
\texttt{kundai.sachikonye@bitspark.com} \\
(+49) 1626386344 \\[6pt]
16 May 2026
\end{flushright}

\vspace{1em}

Prof.\ Dr.\ Lars Lindner \\
Sarcoma / Hyperthermia / Drug Delivery Working Group \\
Medizinische Klinik III (Onkologie) \\
LMU University Hospital Munich

\vspace{1.5em}

\textbf{Re: PhD / Postdoctoral Position --- In Vivo Cancer Therapeutics}

\vspace{1em}

Dear Prof.\ Lindner,

I am applying for the PhD / Postdoctoral position in in vivo cancer therapeutics.
My doctoral research was in Computational Lipidomics at TUM, where I developed
mass spectrometry pipelines specifically for lipid systems characterisation. That
background connects directly to the core analytical problem of your group's
work: thermosensitive liposomal drug delivery depends on phospholipid phase
transitions that are precisely the events that lipidomics mass spectrometry
measures, identifies, and can now predict from first principles.

\textbf{Lipidomics and liposomal drug delivery.}
At TUM / Bayerisches Forschungsinstitut I built the complete data pipeline for
lipid species analysis: peak detection and deconvolution, ion annotation against
lipid databases (LIPID MAPS, LipidBlast), multi-omics statistical modelling,
and federated learning architectures for distributed clinical data. Separately,
at Fraunhofer Institut IGB Stuttgart, I applied dynamic programming and linear
programming methods to the production optimisation of \textit{therapeutic lipids}
--- the same lipid classes that constitute thermosensitive liposomal formulations.
The
\href{https://github.com/fullscreen-triangle/lavoisier}{Lavoisier} framework
I built for mass spectrometry extends this: the virtual instrument at
\href{https://lavoisier.vercel.app/experiment}{lavoisier.vercel.app/experiment}
runs forward simulation of mass spectrometry experiments in the browser, and the
underlying framework processes mzML datasets exceeding 100\,GB, achieving 98.9\,\%
feature extraction accuracy against the NIST reference library.
The phase transition of DPPC/DSPC-based thermosensitive liposomes from gel to
liquid crystalline phase at approximately 42\,$^\circ$C is observable by
mass spectrometry as a characteristic shift in the phospholipid ion population
--- a measurement class the pipeline was built to characterise.

\textbf{Why sarcoma and pancreatic cancer are immunologically cold: a first-principles account.}
My manuscript \textit{On the Thermodynamic Consequences of Categorical Completion
on Biological Systems} models adaptive immune selectivity as a categorical filter
cascade: MHC molecules present peptides preferentially from proteins with high
categorical richness $R$ --- quantifying isoform diversity, conformational entropy,
and post-translational modification landscape. The correlation is $r = 0.73$,
$p < 10^{-12}$ across 2,847 IEDB peptides; adding the $R$ term to NetMHCpan
raises AUC from 0.68 to 0.81.
Sarcoma is paradigmatically collagen-driven: the dense ECM of soft-tissue
sarcomas is dominated by collagen isoforms with $R \approx 10^{5.5}$,
categorically the highest richness in the human proteome, which predicts
them as strong immune targets --- yet sarcomas are immunologically cold.
The model resolves this: high-$R$ target availability does not guarantee
immune activation when the tumour microenvironment suppresses the
MHC-I presentation pathway. Pancreatic cancer presents the same structure
amplified by desmoplastic stroma. The combination of hyperthermia (which
disrupts stromal architecture and temporarily restores MHC-I surface density)
with liposomal drug delivery and immunotherapy is precisely the intervention
the framework predicts will convert these tumours from cold to hot.
I can derive specific predictions about which peptide populations will expand
following hyperthermia treatment, and which drug targets will have the lowest
off-target burden --- predictions testable in your group's in vivo models.

\textbf{Tumour microenvironment as a trajectory completion problem.}
The \href{https://github.com/fullscreen-triangle/syndrome}{Syndrome} framework
computes cellular disease state as a vector $\mathbf{D} \in [0,1]^8$ across
eight oscillator classes. For the hyperthermia / liposomal drug delivery
context, the directly relevant components are the membrane potential oscillator
($D_M$, responding to thermal stress-induced bilayer perturbation), the calcium
oscillator ($D_{Ca}$, which governs the heat-shock response and drives
immunogenic cell death signalling), and the ATP oscillator ($D_{ATP}$, depleted
under the metabolic demands of combined hyperthermia and cytotoxic drug exposure).
The framework computes the trajectory from tumour cell state to the target
state (immunogenic, drug-sensitive, cleared) in $\mathcal{O}(N \log N)$, and
the sparse $\ell_1$ linear programme selects the minimum-intervention combination
--- which of hyperthermia dose, liposomal drug concentration, and immunotherapy
agent --- needed to complete the transition. This is exactly the optimisation
problem your multimodal therapy programme needs to solve.

\textbf{Computational infrastructure for the experimental programme.}
The tools I have deployed address each analytical dimension the programme requires.
\href{https://syndrome-seven.vercel.app}{syndrome-seven.vercel.app} provides
four live immunopeptidology tools operating on the same spectral complementarity
metric $\rho$ (3-channel DFT, 36-dimensional cosine similarity, $O(cL\log L)$):
a neoantigen identifier (dual threshold: $\Delta\rho_\text{self} > 0.15$ for
immune detection, $\rho_\text{MHC} > 0.72$ for MHC presentation) --- directly
applicable to identifying which intracellular peptides become immunogenic after
hyperthermia-mediated lysis; and an antibiotic resistance predictor
(drug--enzyme spectral complementarity $\rho$, validated against known clinical
resistance positions for TEM-1/Ampicillin, GyrA/Ciprofloxacin,
AAC(3)-IV/Gentamicin) --- the same spectral distance applied to chemotherapy
drug--target interactions to predict which mutations would confer resistance
to doxorubicin or other cytotoxic agents.
\href{https://vivid-symbolism.vercel.app/search}{vivid-symbolism.vercel.app/search}
provides sequence homology search underpinning tumour neoantigen identification.
\href{https://remote-photoplethysmography.vercel.app/}{BRUT Observatory} inverts
a layered Beer--Lambert optical model of forehead skin to recover vasodilation
and skin temperature from a standard webcam in real time --- directly applicable
to non-contact thermal monitoring of patients during hyperthermia sessions,
supplementing contact thermometry without additional hardware.
\href{https://shakespear-nine.vercel.app/instrument}{shakespear-nine.vercel.app/instrument}
provides protein folding, trajectory, and dynamics analysis applicable to the
oncoprotein targets released by hyperthermia-mediated cell lysis.
The \href{https://github.com/fullscreen-triangle/borgia}{Borgia} framework
derives molecular spectroscopic properties from bounded phase space geometry ---
the Triple Equivalence Theorem equates oscillatory dynamics, categorical state
enumeration, and partition operations, making it possible to predict the DPPC
and DSPC vibrational spectra and phase transition temperatures from structural
geometry alone, without parameter fitting against empirical data.
The cheminformatics instrument at
\href{https://honjo-masamune.vercel.app/tools}{honjo-masamune.vercel.app/tools}
implements this directly: zero-parameter molecular spectroscopy including docking,
3D structure, and harmonic network analysis, all computed in the browser.

\textbf{Drug-target geometry and therapeutic trajectories.}
The \href{https://crown-prince-four-courners.vercel.app}{Crown Prince} platform
extends the same bounded phase space framework to pharmacology.
Drug-target binding affinity emerges as $K_d = \exp(-\Delta M \cdot \ln b)$
from partition geometry alone --- validated against 39/39 NIST binding measurements
with zero free parameters (50\,MB covering the same pharmacological space as the
12\,GB DrugBank reference). For thermosensitive liposomal delivery, the Therapeutic
Effect Trajectory paper in the platform derives the full PD/PK/ADME profile as
geometric transitions in partition space: drug absorption, distribution, and
clearance as first-principles state trajectories rather than phenomenological
curve-fitting. The Cell Spectral Hologram tool at
\href{https://crown-prince-four-courners.vercel.app/tools/spectrum}{crown-prince-four-courners.vercel.app/tools/spectrum}
accepts a tumour cell type as input and outputs a sparse $\eta^*$ therapeutic
perturbation vector identifying which molecular targets, when modulated, move the
cell from the disease trajectory to the clearance state with minimum cytotoxic
burden --- directly applicable to the multimodal optimisation problem of choosing
hyperthermia dose, liposomal drug concentration, and immunotherapy agent.

\textbf{Background.}
My MSc in Pharmaceutical Biotechnology (HAW Hamburg) covered drug delivery,
formulation science, and pharmacokinetics. My BSc in Biochemistry and Cell
Biology (Jacobs University Bremen) included standard cell culture and molecular
biology laboratory training. My MSc in Computational Biology (Universität
Tübingen) and doctoral research in Computational Lipidomics at TUM gave
me the analytical and statistical depth behind the work described above. I am
not proposing to arrive as an animal model specialist --- that is the training
this position provides. I am proposing that a candidate who can derive from
first principles why sarcoma and pancreatic cancer are immunologically cold,
predict which hyperthermia-driven state transitions will produce immunogenic
cell death, and analyse the resulting lipidomic and proteomic data at scale
will contribute substantively to this programme from the first day.

I am legally resident in Germany with full work authorisation and am available
from October 2026. English is native; German is conversational.

\vspace{1.5em}
Yours sincerely,

\vspace{2.5em}
\textbf{Kundai Farai Sachikonye}

\end{document}
